% Material removed from main.tex during the ergodicity restructuring (October 2026).
% It is kept verbatim as the starting point for a second paper centred on the
% BABP sandwich result (old Corollary 4.2). This file is not compiled on its own;
% labels refer to the pre-restructuring numbering of main.tex.

%==============================================================================
% Old introduction (comparison-principle framing)
%==============================================================================

\section{Introduction}
\label{sec:introduction}

Interacting particle systems are stochastic models in which local random
updates at many sites combine to produce collective behaviour. We study spin
systems, continuous-time Markov processes on~$\cb{0,1}^{\Lambda}$\chg{, where~$\Lambda$ is a countable set of sites,} in which
one spin is updated at a time and the update rule may depend on nearby spins;
see~\cite{Liggett1985} for the general theory. A basic question is how the
local update rule determines the dependence of the process on its initial
condition, in particular its invariant measures and long-time behaviour.

\chg{Comparison principles are among the most useful tools for this
question. By a comparison principle we mean an order on probability laws
that the dynamics preserve. The classical example is monotonicity, where
the order is stochastic domination. A spin system is monotone,
traditionally called attractive, if processes started from configurations
ordered coordinatewise,~$\eta\leq\xi$, can be coupled so that the order is
preserved at all later times~\cite{Harris1977}, \cite[Chapter~III,
Theorem~2.2]{Liggett1985}; equivalently, stochastically ordered initial
laws have stochastically ordered evolutions~\cite[Chapter~II, Theorem~2.2
and Definition~2.3]{Liggett1985}. For monotone spin systems this has the
following standard consequences. Evolutions started from the all-$0$ and
all-$1$ configurations are monotone in time and converge to the lower and
upper invariant measures, every other evolution is sandwiched between them,
and ergodicity reduces to showing that these two measures
coincide~\cite[Chapter~III, Theorem~2.3 and Corollary~2.4]{Liggett1985}.
Moreover, spin systems whose rates are suitably ordered can be compared
with each other~\cite[Chapter~III, Theorem~1.5 and
Corollary~1.7]{Liggett1985}.}

\chg{We start from the simplest spin system, independent Bernoulli refreshes:
every site is resampled from~$\operatorname{Bern}(p)$ at rate one,
regardless of the other spins. These dynamics are monotone, since two
ordered initial configurations driven by the same clocks and the same coins
stay ordered, so they satisfy a comparison principle. We then constrain the dynamics: an update at a site is carried
out only if a local condition on the neighbouring spins is satisfied. This
gives the kinetically constrained spin models
(KCSMs)~\cite{FredricksonAndersen1984,CancriniEtAl2008,HartarskyToninelli2025}.}
For example, in the Fredrickson--Andersen
one-spin facilitated model (FA-1f), an update is allowed when at least one
neighbouring site is vacant. The East model imposes the same type of
constraint but on an oriented neighbourhood graph. Thus these models have a
simple Bernoulli product equilibrium measure, while the kinetic constraint can
make relaxation slow and the out-of-equilibrium behaviour difficult. {\color{red} Natural point to stae open question for ergodicity from unconstrained state for FA. Why to show that this is an avtive research area of interest!} Precise
definitions of the models used in this paper are given in
Section~\ref{sec:setup}.

\chg{The constraint destroys monotonicity, the classical comparison
principle. For example, in FA-1f or East
let~$j\in N(i)$, $\eta(x)=\ind(x\neq i,j)$ and~$\xi(x)=\ind(x\neq i)$.
Then~$\eta\leq\xi$, but an update at~$i$ that refreshes to~$1$ is legal
in~$\eta$ and blocked in~$\xi$, so afterwards~$\eta(i)=1>\xi(i)=0$. The
reason is general. Monotonicity requires the~$0\to1$ rate to increase and
the~$1\to0$ rate to decrease as spins are added~\cite[Chapter~III, Definition~2.1 and
Theorem~2.2]{Liggett1985}, but in a
KCSM both rates are proportional to the same constraint. Hence a KCSM is
monotone only if its constraint is trivial.

This raises the question that motivates the paper: does a comparison
principle survive the kinetic constraint, and does it still give access to
invariant measures and convergence? More generally, can one recover such a
principle for a large class of spin systems whose dynamics are not monotone
in configuration order?}

\chg{We identify a class of spin systems, called \emph{patch-positive}, for which a
comparison principle survives. Patch positivity is an explicit sign
condition on the local update rates, stated in
Definition~\ref{def:patch-positivity} and interpreted probabilistically in
Section~\ref{sec:patches}. The class contains FA-1f and East, but not
FA-$j$f with~$j\geq2$. It also contains models outside the KCSM class, such
as the biased annihilating branching process (BABP) and, after exchanging
the two spin values, the contact process. These examples are worked out in
Section~\ref{sec:applications}.}

\chg{Our first result is a comparison principle for patch-positive spin
systems. It uses an order on probability laws rather than on configurations.} \chg{Each patch-positive system has
a threshold profile~$\mb p^\star=(p_i^\star)_{i\in\Lambda} \in [0,1]^\Lambda$;
for FA-1f and East it equals the refresh density~$p$.} For a finite
set~$A\subseteq\Lambda$, define the centered monomial
\begin{equation*}
\chi_A^\star(\eta)
=
\prod_{i\in A}\bb{\eta(i)-p_i^\star}.
\end{equation*}
\chg{Its expectation~$\mu(\chi_A^\star)$ is a joint moment of the spins
in~$A$, centered at the threshold profile rather than at the mean
of~$\mu$.} For probability measures~$\mu$ and~$\nu$, we write~$\mu\preceq_\star\nu$
when
\begin{equation*}
\mu\bb{\chi_A^\star}
\leq
\nu\bb{\chi_A^\star}
\qquad
\text{for every finite }A\subseteq\Lambda.
\end{equation*}
We call this the \emph{centered monomial order}. \chg{Writing
$\eta(i)=(\eta(i)-p_i^\star)+p_i^\star$ and expanding the product shows
that~$\mu\preceq_\star\nu$ implies
\begin{equation*}
\mu\bb{\eta(i)=1\text{ for all }i\in A}
\leq
\nu\bb{\eta(i)=1\text{ for all }i\in A}
\qquad
\text{for every finite }A\subseteq\Lambda,
\end{equation*}
so every finite set is at least as likely to be fully occupied under~$\nu$
as under~$\mu$. The centered monomial order strengthens this comparison. It
is easiest to read for laws with~$\mu(\chi_A^\star)\geq0$ for all~$A$,
which is where our comparison results apply. {\color{blue}\textbf{reformulate, actually improve and give the simplest example of a class of probabitlity measures that satisfy a centered monomial order...}} The Bernoulli product law with
profile~$\mb p^\star$ is the smallest such law, since all its nonconstant
centered monomials have expectation zero, and~$\nu$ lies above~$\mu$ if on
every finite set the spins exceed~$\mb p^\star$ jointly at least as
strongly under~$\nu$ as under~$\mu$.} \chg{The result of
Theorem~\ref{thm:centered-monomial-monotonicity} is that patch-positive
dynamics preserve this order: writing~$P_t$ for the Markov semigroup,}
\begin{equation*}
\mu\preceq_\star\nu
\quad\Longrightarrow\quad
\mu P_t\preceq_\star\nu P_t
\qquad
\text{for every }t\geq0.
\end{equation*}
\chg{Thus the configuration
order is lost along trajectories, but a comparison principle survives at the
level of probability laws. Patch positivity is sufficient for this
comparison principle, but not necessary: the semigroup can preserve the
centered monomial order also for spin systems that are not patch-positive.}

\chg{This comparison principle has the same kinds of consequences as in the
  monotone case.} {\color{blue}\textbf{weitermachen}}

{\color{red}
\begin{itemize}
\item Now explain analog to monotonicity what the consequences of this comparison principle are:
  \begin{itemize}
  \item Essentially all carry over just not on the full set but on a constraint set of initial laws $\mathcal{M}_{*}$.
\item Need to explain what is~$\mathcal{M}_{*}$. It is the set of highly contrained measures. as it is defined as the measures with nonnegative expectation of centered monomials. 

\item Sandwiching works on ~$\mathcal{M}_{*}$ (refer to above Ligget....and Corollary 3.6), application is: For BABP, known exponential convergence from homogeneous
product measures can be used as the endpoints for a sandwiching argument. This gives
exponential convergence to equilibrium for a class of dependent and non-translation-invariant
initial laws; see Corollary 4.2.
\item comparison principle comparing evolution of different spin system also works in the following way......(Corollary 3.7))
\end{itemize}
\item Observation: Conseqeunces of monomial order comparison principle are weaker than the consequences of monotonicity.
  \begin{itemize}


  \item Example of type 1:  
  
  
  \end{itemize}
\end{itemize}
}







{\color{red} Patch representation allows to compare contributions patches of different models this yields beautful results:
\begin{itemize}
  \item On neighbourhood graphs of
polynomial growth, if a patch-positive system contains pure deaths at a uniformly positive
rate, then a broad class (i.e.~$\mathcal{M}_{-}$) of initial laws converges to a common invariant measure. This class
includes every configuration with only finitely many 0s. See Theorem 3.9.
\item Consequences of Theorem 3.9:
  \begin{itemize}
  \item   ~$p^{*} \leq \frac{1}{2}$ (meaning a lot of activity)) then~$\mathcal{M}_{-}$ is everything. Corollary 3.10.
  \item  Patch positive, polynomial graph plus oriented interaction like East model are uniformly ergodic. (not exponential). Corollary 3.11. Good result! Contains the perturbed east model (Corollary 4.6). 
  \end{itemize}
\end{itemize}
  




}

The Bernoulli product law
with profile~$\mb p^\star$ evolves monotonically upward in time and converges to an
invariant measure. When~$\mb p^\star\leq\mbs{\tfrac12}$ the same is true about the evolution of the all~$1$ configuration.  More generally, the order admits a sandwiching argument, see Corollary~\ref{cor:extremal-limits-comparison}. The order also
compares different dynamics. In particular, if a patch-positive system is
obtained from another one by adding independent~$1\to0$ transitions, which we call pure deaths, then
for the initial laws covered by Corollary~\ref{cor:pure-death-comparison-revision},
the expectations of all finite spin monomials are smaller in the system with
the additional transitions.

The comparison principle also leads to convergence results. On neighbourhood
graphs of polynomial growth, if a patch-positive system contains pure deaths at a uniformly positive rate, then a broad class of
initial laws converges to a common invariant measure. This class includes every
configuration with only finitely many~$0$s. See
Theorem~\ref{thm:common-invariant-limit}.

There are two routes from this result to convergence uniformly
over all initial configurations. If
$\mb p^\star\leq\mbs{\tfrac12}$, the system is uniformly exponentially
ergodic; see Corollary~\ref{cor:uniform-exponential-ergodicity}. Alternatively,
if the neighbourhood graph is orientable, then the system is uniformly
ergodic without any restriction of this kind on~$\mb p^\star$; see
Corollary~\ref{cor:oriented-ergodicity}. The perturbed East model gives the main example
of the second route; see Corollary~\ref{cor:perturbed-east}.

The applications also show how the centered monomial order can be combined
with information specific to a model. For BABP, known exponential convergence
from homogeneous product measures can be used as the endpoints for a sandwiching argument. This gives exponential convergence to equilibrium for a
class of dependent and non-translation-invariant initial laws; see
Corollary~\ref{cor:babp-centered-convergence}. For one-dimensional, one-sided nearest neighbor spin systems, called Simple~\cite{GluchowskiMenz2025,GluchowskiMenz2026}, patch positivity together with orientability gives a new sufficient
condition for uniform ergodicity under positive rates, extending the regions
covered by earlier criteria; see
Corollary~\ref{cor:simple-positive-rates}.

The main construction is a representation of the action of the Markov semigroup on monomials. For~$\chi_A(\eta)=\prod_{i\in A}\eta(i)$, the starting point is
the set-valued Feynman--Kac representation
\begin{equation}
\label{eq:intro-signed-duality}
P_t\chi_A(\eta)
=
\E_A\Cb{
\sigma_t
\exp\bb{\int_0^t V(A_s)\,ds}
\chi_{A_t}(\eta)
}.
\end{equation}
Here~$A_s$ is the finite set-valued dual process,~$\sigma_s\in\cb{+,-}$, and~$V$
is the Feynman--Kac potential. Set-valued dual constructions for spin systems
are standard; see, for example,
\cite{Harris1978,Griffeath1979,Gray1986,JansenKurt2014}. The important feature
for us is that the right-hand side of~\eqref{eq:intro-signed-duality} is
signed, which is an obstruction for monotonicity or convergence arguments.

The new step is to reveal only part of the graphical construction. We define a~$\sigma$-algebra
$\mathcal G_t$ that records interactions between sites that occur in the
dual up to time~$t$, including their times and the sites they connect, while
forgetting their type and the purely local events between them. Once this
information is fixed, the recorded interactions divide each site line~$\cb{i}\times[0,t]$ into
consecutive time intervals. We call these intervals \emph{patches}. Thus a
patch is a piece of one site line between two consecutive recorded
interactions, or between such an interaction and the observation time.

Conditioning on~$\mathcal G_t$ makes the reason for this choice visible. The
forgotten randomness has a conditional law that factorizes over the patches,
and the signed random variable in~\eqref{eq:intro-signed-duality} is itself a
product of local factors. Consequently, schematically,
\begin{equation}
\label{eq:intro-patch-factorization}
\E_A\Cb{
\sigma_t
\exp\bb{\int_0^t V(A_s)\,ds}
\chi_{A_t}(\eta)
\,\middle|\,
\mathcal G_t
}
=
\prod_{\substack{ \text{patch }P\text{ ending}\\\text{before }t}} C(P)
\prod_{\substack{\text{patch }P\text{ reaching}\\\text{time }t}}
C\bb{\eta\bb{i(P)},P}.
\end{equation}
The notation in~\eqref{eq:intro-patch-factorization} is only schematic. 
Patches and their contributions are defined precisely in
Section~\ref{sec:patches}. The point is this conditioning allows for local averaging, which gives explicit contributions and captures cancellations from sign process in~\eqref{eq:intro-signed-duality}.

The factorization~\eqref{eq:intro-patch-factorization} also explains the
condition of patch positivity. A patch that ends before time~$t$ leaves no
spin variable in the representation and, after the local averaging above,
contributes a number~$C(P)$. Patch positivity means that every such
contribution is nonnegative. A patch that reaches time~$t$ instead contributes
$C(z,P)$, an affine function of the remaining spin argument~$z$. For a site
$i$, the threshold~$p_i^\star$ is the smallest~$p\in[0,1]$ for which
\begin{equation*}
C(p,P)\geq0
\end{equation*}
for every patch reaching time~$t$ on the site~$i$. Thus the threshold profile
is the minimal density profile~$\mb p$ such that for a product measure~$\mu = \otimes_{i\in\Lambda}\operatorname{Ber}(p_i)$ the~$\mu P_t \chi_A$ can be represented using only nonnegative factors. This supplies enough monotonicity to yield the order preservation result and earlier stated convergence results.

The duality in~\eqref{eq:intro-signed-duality} belongs to the standard
graphical and set-valued duality framework for spin systems
\cite{Harris1974,Harris1978,Griffeath1979,HolleyStroock1979,Gray1986,JansenKurt2014}.
The new ingredient here is the choice of the smaller~$\sigma$-algebra~$\mathcal G_t$.
It is this conditional factorization, rather than the duality itself, that
exposes the local cancellations used in our arguments. Information
percolation also records updates that carry information
\cite{LubetzkySly2016}, but uses them for a different decomposition.

Classical comparison results for monotone spin systems are based on couplings
that preserve the coordinatewise order on configurations and on related
association inequalities
\cite{Harris1977,Liggett1985,Liggett2006}. Duality can also yield correlation
inequalities in particular interacting particle systems
\cite{GiardinaRedigVafayi2010}. Patch-positive systems need not admit such a
coupling. Instead, the local sign conditions revealed by the patch
factorization lead to the centered monomial order on probability laws. Thus
the comparison principle developed here plays a similar role for the models
considered in this paper, while being formulated in terms of centered
monomials rather than configuration order.

The paper is organized so that the results and applications can be read
without the patch construction. Section~\ref{sec:setup} introduces the general
spin-system framework and the examples used throughout the paper.
Section~\ref{sec:main-results} defines patch positivity and the centered
monomial order and states the comparison and convergence results.
Section~\ref{sec:applications} applies them to FA-1f, the East model, BABP,
the contact process, and simple spin systems. The remaining sections develop
the proofs. Section~\ref{sec:dual} constructs the signed dual and its graphical
representation, and Section~\ref{sec:patches} introduces the conditioning,
patches, and patch factorization described above. Section~\ref{sec:proofs}
uses this representation to prove the main results.
Section~\ref{sec:discussion} discusses the scope and limitations of the
approach. The appendices contain technical details of main proofs.

%==============================================================================
% Section 2: contact process example
%==============================================================================

\begin{example}[Contact process]
The contact process~\cite{Harris1974} is usually written with state~$1$
infected and state~$0$ healthy. With recovery rate one and infection
parameter~$\beta\geq0$, it is given by
\begin{equation*}
\lambda_i(\eta)
=
1+\beta\sum_{j\in N(i)}\eta(j),
\qquad
\rho_i(\eta)
=
\frac{\beta\sum_{j\in N(i)}\eta(j)}{1+\beta\sum_{j\in N(i)}\eta(j)}.
\end{equation*}
\end{example}

%==============================================================================
% Section 3: centered monomial order, Theorem 3.5, Corollaries 3.6 and 3.7
%==============================================================================

\begin{definition}[Centered monomial order]
For~$A\Subset\Lambda$, define the centered monomial
\begin{equation*}
\chi_A^\star(\eta)
=
\prod_{i\in A}\bb{\eta(i)-p_i^\star},
\qquad
\chi_\vn^\star=1.
\end{equation*}
For probability measures~$\mu$ and~$\nu$ on~$\Omega$, we set
\begin{equation}
\label{eq:centered-monomial-order}
\mu\preceq_\star\nu \qquad \iff \quad \mu\bb{\chi_A^\star}
\leq
\nu\bb{\chi_A^\star}
\qquad
\text{for every }A\Subset\Lambda.
\end{equation}
We also define
\begin{equation}
\label{eq:m-star-definition-revision}
\mathcal M_\star
=
\cb{
\mu:
\mu\bb{\chi_A^\star}\geq0
\text{ for every }A\Subset\Lambda
}.
\end{equation}
\end{definition}

To understand this order it is useful to see how it compares the product measures. If~$\mb p^\star\leq\mb p\leq\mb q$, then~$\mu_{\mb p^*}\preceq_* \mu_{\mb p}\preceq_* \mu_{\mb q}$. Product measures with profiles that stray below~$\mb p^*$ are not comparable in the order.


\begin{remark}
The measure~$\mu_{\mb p^\star}$ is the minimum of~$\mathcal M_\star$ under
$\preceq_\star$, since all of its nonconstant centered monomials have
expectation zero. In general~$\mathcal M_\star$ need not have a maximum. If~$\mb p^\star\leq\mbs{\tfrac12}$ then~$\mu_{\mb 1}$ is a maximum.
\end{remark}

\begin{theorem}[Centered monomial monotonicity]
\label{thm:centered-monomial-monotonicity}
Suppose that the spin system is patch-positive. If~$\mu$ and~$\nu$ are signed measures on~$\Omega$ and~$\mu\preceq_\star\nu$, then~$\mu P_t\preceq_\star\nu P_t$ for every~$t\geq 0$.
Moreover,~$P_t$ preserves~$\mathcal M_\star$. The centered monomial order also
implies the corresponding order for ordinary monomials: if
$\mu\preceq_\star\nu$, then
\begin{equation}
\label{eq:ordinary-monomial-comparison}
(\mu P_t)(\chi_A)
\leq
(\nu P_t)(\chi_A)
\qquad
\text{for every }A\Subset\Lambda\text{ and }t\geq0.
\end{equation}
\end{theorem}

The theorem gives monotone evolution at the bottom of
$\mathcal M_\star$, and also at the top whenever a maximum exists.

\begin{corollary}[Extremal limits and sandwiching]
\label{cor:extremal-limits-comparison}
Suppose that the spin system is patch-positive.
\begin{enumerate}[label=\textup{(\roman*)}]
\item The sequence~$\mu_{\mb p^\star}P_t$ is increasing in~$t$ under
$\preceq_\star$ and converges weakly to an invariant measure~$\pi^-\in\mathcal M_\star$.

\item If~$\mb p^\star\leq\mbs{\tfrac12}$, then
$\mu_{\mb 1}P_t$ is decreasing in~$t$ under~$\preceq_\star$ and converges
weakly to an invariant measure~$\pi^+\in\mathcal M_\star$.

\item More generally, suppose that
\begin{equation*}
\underline\mu\preceq_\star\mu\preceq_\star\overline\mu
\end{equation*}
and that both~$\underline\mu P_t$ and~$\overline\mu P_t$ converge weakly to
the same measure~$\pi$. Then~$\mu P_t\Rightarrow\pi$.
\end{enumerate}
\end{corollary}

\begin{corollary}[Pure-death comparison]
\label{cor:pure-death-comparison-revision}
Let~$\Lcal$ and~$\Lcal'$ be spin-system generators. Suppose that~$\Lcal$ is
patch-positive and that
\begin{equation}
\label{eq:pure-death-decomposition-revision}
\Lcal f(\eta)
=
\Lcal'f(\eta)
+
\sum_{i\in\Lambda}d_i
\bb{f\bb{\eta^{i,0}}-f(\eta)},
\qquad d_i\geq0,
\end{equation}
so~$\mathcal{L}$ is perturbation of~$\mathcal{L}'$ by adding independent deaths.
Then~$\Lcal'$ is patch-positive with the same threshold profile~$\mb p^\star$.
If~$P_t'$ denotes its semigroup, then
\begin{equation}
\label{eq:pure-death-comparison-revision}
(\mu P_t)(\chi_A)
\leq
(\mu P_t')(\chi_A)
\end{equation}
for every~$\mu\in\mathcal M_\star$, every~$A\Subset\Lambda$, and every
$t\geq0$.
\end{corollary}

Thus, within~$\mathcal M_\star$, adding deaths can only decrease the expectations of the monomials. The result becomes stronger if the invariant measures for one of the systems are known, for example for perturbed KCSMs.

%==============================================================================
% Section 4.2-4.3: BABP sandwich (Corollary 4.2) and contact process
%==============================================================================

Qualitative convergence of BABP from translation-invariant initial laws was
proved by Sudbury~\cite{Sudbury1997}, while quasi-duality and thinning
relations were developed by Lloyd and Sudbury~\cite{LloydSudbury1997}.
On~$\Z^d$, every homogeneous product law~$\mu_u$ with~$u<1$ converges
exponentially to~$\mu_p$
\cite[Application~2 and Section~5.1.1]{MartinelliShapiraToninelli2025}.
The centered monomial order allows these product laws to be used as comparison
laws for more general initial distributions.

\begin{corollary}[Centered monomial comparison for BABP]
\label{cor:babp-centered-convergence}
Consider nearest-neighbour BABP on~$\Z^d$. Let~$0\leq u\leq v<1$ satisfy
\begin{equation*}
v\geq p,
\qquad
\abs{u-p}\leq v-p.
\end{equation*}
If a probability measure~$\nu$ satisfies
\begin{equation}
\label{eq:babp-centered-sandwich}
(u-p)^{\abs A}
\leq
\nu\bb{\prod_{i\in A}\bb{\eta(i)-p}}
\leq
(v-p)^{\abs A}
\qquad
\text{for every }A\Subset\Z^d,
\end{equation}
then~$\nu P_t$ converges exponentially to~$\mu_p$ on local functions.
\end{corollary}

\begin{proof}
Condition~\eqref{eq:babp-centered-sandwich} is exactly
\begin{equation*}
\mu_u\preceq_\star\nu\preceq_\star\mu_v.
\end{equation*}
By Theorem~\ref{thm:centered-monomial-monotonicity}, this ordering is preserved
by~$P_t$. Both endpoint laws converge exponentially to~$\mu_p$, so every
centered monomial expectation under~$\nu P_t$ is squeezed between two
quantities converging exponentially to its~$\mu_p$ expectation. Since centered
monomials form a basis for the functions on every finite set of sites, the same
holds for every local function.
\end{proof}

This includes dependent and non-translation-invariant initial laws. In
particular, taking~$u=p$, every inhomogeneous product law with one-site
parameters
\begin{equation*}
p\leq p_i\leq v<1
\end{equation*}
satisfies~\eqref{eq:babp-centered-sandwich} and therefore converges
exponentially to~$\mu_p$. This allows parameters above the bound
$p_i\leq\sqrt p$ appearing in
\cite[Remark~5.5]{MartinelliShapiraToninelli2025}, provided they are uniformly
bounded away from~$1$.

\subsection{Contact process}
\label{subsec:contact-process-applications}

In the usual labeling from Section~\ref{sec:setup}, the contact process has
state~$1$ infected and
\begin{equation*}
\lambda_i(\eta)
=
1+\beta\sum_{j\in N(i)}\eta(j).
\end{equation*}
Its nonconstant coefficients are positive, so this labeling does not satisfy
the first condition in~\eqref{eq:patch-positivity}. Exchange the roles of
states~$0$ and~$1$. In the resulting variables,
\begin{equation*}
\lambda_i(\eta)
=
1+\beta\sum_{j\in N(i)}\bb{1-\eta(j)},
\qquad
\rho_i(\eta)
=
\frac{1}{1+\beta\sum_{j\in N(i)}\bb{1-\eta(j)}}.
\end{equation*}
Thus~$\lambda_i\rho_i=1$, while the only nonzero nonconstant coefficients of
$\lambda_i$ are
\begin{equation*}
\widehat{\lambda_i}\bb{\cb{j}}=-\beta,
\qquad j\in N(i).
\end{equation*}
For every nonempty~$S\subseteq N(i)$,
\begin{equation*}
\widehat{\lambda_i\bb{\rho_i-\rho_i(\mb0)}}(S)
=
-\rho_i(\mb0)\widehat{\lambda_i}(S)
\geq0.
\end{equation*}
Hence the complemented contact process is patch-positive. Since
$\lambda_i\rho_i$ is constant, its threshold profile is~$\mb p^\star=\mb0$.

The contact process is already monotone in the usual configuration order, and
the classical monotone coupling gives stronger comparisons than the centered
monomial order. We only mention it to demonstrate the scope of patch
positivity.

%==============================================================================
% Section 4.4: comparison with the Bernoulli measure (Corollary 4.4)
%==============================================================================

The limiting measure can also be compared with the Bernoulli product measure
of the unperturbed system.

\begin{corollary}[Comparison with the Bernoulli measure]
\label{cor:perturbed-fa-babp-comparison}
Under the assumptions of
Proposition~\ref{prop:perturbed-fa-babp-convergence},
\begin{equation}
\label{eq:perturbed-fa-babp-stationary-comparison}
p^{\abs A}
\leq
\pi(\chi_A)
\qquad
\text{for every }A\Subset\Lambda.
\end{equation}
If~$\mb s=\mb p$, then~$\pi=\mu_p$.
\end{corollary}

\begin{proof}
Set
\begin{equation*}
d_i=\varepsilon\frac{s_i-p}{p}\geq0
\end{equation*}
and define
\begin{equation*}
\widetilde\Lcal
=
\Lcal^{\varepsilon,\mb s}
+
\sum_{i\in\Lambda}d_i\bb{E_{i,0}-I}.
\end{equation*}
Then
\begin{equation*}
\widetilde\Lcal
=
\Lcal
+
\sum_{i\in\Lambda}
\frac{\varepsilon s_i}{p}\bb{E_{i,p}-I}.
\end{equation*}
Hence~$\widetilde\Lcal$ is patch-positive with threshold~$p$, and~$\mu_p$
is invariant for its semigroup~$\widetilde P_t$. By
Corollary~\ref{cor:pure-death-comparison-revision},
\begin{equation*}
p^{\abs A}
=
(\mu_p\widetilde P_t)(\chi_A)
\leq
(\mu_pP_t^{\varepsilon,\mb s})(\chi_A).
\end{equation*}
Letting~$t\to\infty$ gives~\eqref{eq:perturbed-fa-babp-stationary-comparison}.
If~$\mb s=\mb p$, then~$\mu_p$ is itself invariant for
$\Lcal^{\varepsilon,\mb p}$, so convergence from~$\mu_p$ gives
$\pi=\mu_p$.
\end{proof}

%==============================================================================
% Section 7.2: proofs of Theorem 3.5 and Corollary 3.6
%==============================================================================

\subsection{Centered monomial monotonicity}
\label{subsec:proof-centered-monomial-monotonicity}

The proof of Theorem~\ref{thm:centered-monomial-monotonicity} is based on the
action of the generator in the centered-monomial basis. Patch positivity and
the threshold bound from Proposition~\ref{prop:patch-contribution-positivity}
make every off-diagonal coefficient in this basis nonnegative. The resulting semigroup therefore preserve nonnegative centered moments.

\begin{proof}[Proof of Theorem~\ref{thm:centered-monomial-monotonicity}]
If~$i\in A$, then
\begin{equation*}
\lambda_i\bb{E_{i,\rho_i}-I}\chi_A^\star
=
\lambda_i\bb{\rho_i-p_i^\star}\chi_{A\setminus\cb{i}}^\star
-
\lambda_i\chi_A^\star.
\end{equation*}
Expanding~$\lambda_i\rho_i$ and~$\lambda_i$ in~$\mathcal{L}\chi_A^*$ gives
\begin{equation}
\label{eq:centered-monomial-generator-expansion}
\Lcal\chi_A^\star
=
\sum_{i\in A}\Cb{
-\lambda_i(\mb 0)\chi_A^\star
+
\sum_{S\subseteq N(i)}
\bb{
\widehat{\lambda_i\rho_i}(S)
-
p_i^\star\widehat{\lambda_i}(S)
}
\chi_S\chi_{A\setminus\cb{i}}^\star
+
\sum_{\substack{S\subseteq N(i)\\S\neq\vn}}
\bb{-\widehat{\lambda_i}(S)}
\chi_S\chi_A^\star
}.
\end{equation}

Patch positivity implies~$-\widehat{\lambda_i}(S)\geq0$ for every nonempty~$S\subseteq N(i)$.
The other coefficients in~\eqref{eq:centered-monomial-generator-expansion}
are also nonnegative:
\begin{equation}
\label{eq:centered-generator-positive-coefficients}
\widehat{\lambda_i\rho_i}(S)
-
p_i^\star\widehat{\lambda_i}(S)
\geq0
\qquad
\text{for every }S\subseteq N(i).
\end{equation}
For~$S=\vn$, this is
\begin{equation*}
\lambda_i(\mb 0)
\bb{\rho_i(\mb 0)-p_i^\star}
\geq0
\end{equation*}
by Proposition~\ref{prop:patch-contribution-positivity}. If~$S\neq\vn$ and
$\widehat{\lambda_i}(S)<0$, inequality
\eqref{eq:centered-generator-positive-coefficients} follows directly from the
definition~\eqref{eq:threshold-profile} of~$p_i^\star$. If
$\widehat{\lambda_i}(S)=0$, patch positivity gives
\begin{equation*}
\widehat{\lambda_i\rho_i}(S)
=
\widehat{\lambda_i\bb{\rho_i-\rho_i(\mb 0)}}(S)
\geq0.
\end{equation*}

It remains to verify that a product of a monomial and a centered monomial has nonnegative coefficient expansion in the centered monomial basis. For finite~$S,B\subseteq\Lambda$, we have
\begin{equation}
\label{eq:ordinary-centered-product-expansion}
\chi_S\chi_B^\star
=
\bb{\prod_{j\in S\cap B}\bb{1-p_j^\star}}
\sum_{R\subseteq S}
\bb{\prod_{j\in S\setminus R}p_j^\star}
\chi_{R\cup(B\setminus S)}^\star.
\end{equation}
Indeed, on~$S\cap B$ we use
\begin{equation*}
\eta(j)\bb{\eta(j)-p_j^\star}
=
\bb{1-p_j^\star}\eta(j),
\end{equation*}
and then expand each remaining factor as~$\eta(j)=p_j^\star+\bb{\eta(j)-p_j^\star}$. Since~$0\leq p_j^\star\leq1$, every coefficient in
\eqref{eq:ordinary-centered-product-expansion} is nonnegative. Combining this
with~\eqref{eq:centered-monomial-generator-expansion} and
\eqref{eq:centered-generator-positive-coefficients} gives
\begin{equation}
\label{eq:centered-generator-metzler}
\Lcal\chi_A^\star
=
\sum_{B\Subset\Lambda}M(A,B)\chi_B^\star,
\qquad
M(A,B)\geq0
\quad\text{for }B\neq A.
\end{equation}
For each~$A$, only the coefficients~$M(A,B)$ for~$B\subseteq N_*(A)$ are nonzero, and the matrix~$M$ is Metzler, so for each~$t\geq0$ the matrix~$e^{tM}=\sum_{n\in\N}\frac{t^n}{n!}M^n$ is well-defined and has nonnegative entries; see, for example,~\cite{ArendtEtAl1986}. If~$\sigma$ is a
finite signed measure whose centered moments are all nonnegative, then
\begin{equation*}
(\sigma P_t)(\chi_A^\star)
=
\sum_{B\Subset \Lambda}
e^{tM}(A,B)\,\sigma(\chi_B^\star)
\geq0.
\end{equation*}
Taking~$\sigma=\nu-\mu$ proves
$\mu P_t\preceq_\star\nu P_t$ whenever~$\mu\preceq_\star\nu$. The same argument with~$\nu\in\mathcal M_\star$ and~$\mu = 0$ shows that~$P_t$ preserves~$\mathcal M_\star$.

 To obtain the comparison for
monomials, expand
\begin{equation*}
\chi_A
=
\sum_{B\subseteq A}
\bb{\prod_{i\in A\setminus B}p_i^\star}
\chi_B^\star.
\end{equation*}
All coefficients are nonnegative. Therefore
$\mu P_t\preceq_\star\nu P_t$ implies
\begin{equation*}
(\mu P_t)(\chi_A)
\leq
(\nu P_t)(\chi_A)
\end{equation*}
for every~$A\Subset\Lambda$, which proves
Theorem~\ref{thm:centered-monomial-monotonicity}.
\end{proof}

\begin{proof}[Proof of Corollary~\ref{cor:extremal-limits-comparison}]
Since~$\mu_{\mb p^\star}$ is the minimal element of~$\mathcal M_\star$ and~$P_t$ preserves~$\mathcal M_\star$, for every~$s\geq0$,
\begin{equation*}
\mu_{\mb p^\star}
\preceq_\star
\mu_{\mb p^\star}P_s.
\end{equation*}
Applying Theorem~\ref{thm:centered-monomial-monotonicity} and using the
semigroup property gives
\begin{equation*}
\mu_{\mb p^\star}P_t
\preceq_\star
\mu_{\mb p^\star}P_{t+s}
\qquad
\text{for every }s,t\geq0.
\end{equation*}
Thus, for every~$A\Subset\Lambda$, the centered moment~$(\mu_{\mb p^\star}P_t)(\chi_A^\star)$ is increasing in~$t$ and bounded,
so it has a limit. Any subsequential weak limit of~$\mu_{\mb p^\star}P_t$ therefore has the same centered monomial moments.
Since centered monomials form a basis for the functions of every finite set of
spins, these moments determine the weak limit. Hence~$\mu_{\mb p^\star}P_t\Rightarrow\pi^-$ for some probability measure~$\pi^-$. The cone~$\mathcal M_\star$ is closed under weak convergence, so~$\pi^-\in\mathcal M_\star$. Finally, for every~$s\geq0$, the Feller
property and the semigroup identity
\begin{equation*}
\mu_{\mb p^\star}P_{t+s}
=
\bb{\mu_{\mb p^\star}P_t}P_s
\end{equation*}
give~$\pi^-P_s=\pi^-$.

Suppose now that~$\mb p^\star\leq\mbs{\tfrac12}$. Then~$\mu_{\mb1}$ is
the maximal element of~$\mathcal M_\star$, and the same
argument as above gives~$\mu_{\mb1}P_t\Rightarrow\pi^+$ for some invariant~$\pi^+\in\mathcal M_\star$.

For part~\textup{(iii)}, suppose that~$\underline\mu\preceq_\star\mu\preceq_\star\overline\mu$. By
Theorem~\ref{thm:centered-monomial-monotonicity},
\begin{equation*}
\underline\mu P_t
\preceq_\star
\mu P_t
\preceq_\star
\overline\mu P_t
\qquad
\text{for every }t\geq0.
\end{equation*}
Hence, for every~$A\Subset\Lambda$,
\begin{equation*}
(\underline\mu P_t)(\chi_A^\star)
\leq
(\mu P_t)(\chi_A^\star)
\leq
(\overline\mu P_t)(\chi_A^\star).
\end{equation*}
If the two outer measures converge weakly to~$\pi$, the outer terms both
converge to~$\pi(\chi_A^\star)$, so the middle term does as well. Since the
centered monomials form a basis for the functions of every finite set of spins,
expectations of every local function converge to their~$\pi$-expectation.
Thus~$\mu P_t\Rightarrow\pi$.
\end{proof}

%==============================================================================
% Section 7.3: original proof of Corollary 3.7 (pure-death comparison)
%==============================================================================

\subsection{Pure-death comparison}
\label{subsec:proof-pure-death-comparison}

\begin{proof}[Proof of Corollary~\ref{cor:pure-death-comparison-revision}]
Let~$H((A,\sigma),\eta)=\sigma\chi_A(\eta)$ be the duality function from
Section~\ref{sec:dual}. The pure-death term in
\eqref{eq:pure-death-decomposition-revision} acts diagonally on~$H$:
\begin{equation*}
\sum_{i\in\Lambda}d_i\bb{E_{i,0}-I}H((A,\sigma),\eta)
=
-\bb{\sum_{i\in A}d_i}H((A,\sigma),\eta).
\end{equation*}
Hence, if
\begin{equation*}
\Lcal H
=
\Dcal H+V(A)H,
\end{equation*}
then
\begin{equation*}
\Lcal'H
=
\Dcal H
+
\bb{V(A)+\sum_{i\in A}d_i}H.
\end{equation*}
Thus the two systems have the same signed dual generator and therefore the same
dual process, successful-interaction skeleton, patch families, and
patch laws. Their potentials satisfy
\begin{equation}
\label{eq:pure-death-potential-comparison}
V_i'=V_i+d_i.
\end{equation}

We next compare the patch contributions between the two systems. Writing~$\psi_i'$ for the function associated with~$\Lcal'$, we
have, by~\eqref{eq:dual-to-spin-contribution-data},
\begin{align*}
\psi_i(\Delta,z)
&=
\delta_i(\varnothing)
\int_0^\Delta
e^{-\bb{\lambda_i'(\mb 0)+d_i}u}\,du
+
z e^{-\bb{\lambda_i'(\mb 0)+d_i}\Delta},
\\
\psi_i'(\Delta,z)
&=
\delta_i(\varnothing)
\int_0^\Delta
e^{-\lambda_i'(\mb 0)u}\,du
+
z e^{-\lambda_i'(\mb 0)\Delta}.
\end{align*}
Hence
\begin{equation}
\label{eq:pure-death-psi-comparison}
0\leq\psi_i(\Delta,z)\leq\psi_i'(\Delta,z)
\qquad
\text{for }\Delta\in[0,\infty]\text{ and }z\in[0,1].
\end{equation}

First consider a bulk patch with incoming initial interaction, i.e. of type~$\mathsf{II}$ or~$\mathsf{IO}$. The first two
rows of~\eqref{eq:full-contribution-formulas},
together with~\eqref{eq:pure-death-potential-comparison} and
\eqref{eq:pure-death-psi-comparison}, give
\begin{equation*}
0\leq C(P)\leq C'(P).
\end{equation*}
If a bulk patch starts with an outgoing interaction, i.e. is of type~$\mathsf{OI}$ or~$\mathsf{OO}$, and the interaction has target~$S$, then patch positivity gives
\begin{equation*}
\beta_i(S)\sigma_i^\beta(S)
=
-\widehat{\lambda_i}(S)
\geq0.
\end{equation*}
The remaining rows of~\eqref{eq:full-contribution-formulas} therefore give the
same inequality. We have proved
\begin{equation}
\label{eq:pure-death-bulk-patch-comparison}
0\leq C(P)\leq C'(P)
\qquad
\text{for every bulk patch }P.
\end{equation}
Moreover, if~$P$ ends with an outgoing interaction, i.e. is of type~$\mathsf{IO}$ or~$\mathsf{OO}$, then its contribution is~$e^{V_i\Delta}$, where
$i=i(P)$ and~$\Delta=e(P)-s(P)$. Thus
\begin{equation}
\label{eq:pure-death-outgoing-terminal-factor}
C(P)
=
e^{-d_{i(P)}\bb{e(P)-s(P)}}C'(P).
\end{equation}

Now we move on to comparing end patch contributions. By Proposition~\ref{prop:patch-contribution-positivity},
\eqref{eq:pure-death-bulk-patch-comparison} implies that~$\Lcal'$ is
patch-positive and its threshold profile is the same as that of~$\Lcal$. Indeed,~$\lambda_i'(S) = \lambda_i(S)$ for~$S\neq \vn$, while~$\lambda_i\rho_i=\lambda_i'\rho_i'$,
and definition of~$\mb p^*$ in~\eqref{eq:threshold-profile} only involves coefficients for~$S\neq \vn$. For an end patch~$P$ based at
site~$i$, set
\begin{equation*}
\kappa(P)=\partial_zC(z,P),
\qquad
\kappa'(P)=\partial_zC'(z,P).
\end{equation*}
Since~$\varphi_i$ is unchanged, the formulas in
\eqref{eq:end-contribution-formulas} give
\begin{equation}
\label{eq:pure-death-end-patch-comparison}
0\leq\kappa(P)\leq\kappa'(P),
\qquad
0\leq C\bb{p_i^\star,P}
\leq
C'\bb{p_i^\star,P}.
\end{equation}
The second inequality also uses Proposition~\ref{prop:patch-contribution-positivity}
for the two patch-positive systems. Since~$C(z,P)$ is affine in~$z$,
\begin{equation*}
C\bb{\eta(i),P}
=
C\bb{p_i^\star,P}
+
\kappa(P)\bb{\eta(i)-p_i^\star}.
\end{equation*}
Distinct end patches are based at distinct sites. Therefore
\begin{equation}
\mu\bb{
\prod_{P\in\mathcal E_t}
C\bb{\eta\bb{i(P)},P}
} =
\sum_{\mathcal Q\subseteq\mathcal E_t}
\mu\bb{
\chi_{\cb{i(P):P\in\mathcal Q}}^\star
}
\prod_{P\in\mathcal Q}\kappa(P)
\prod_{P\in\mathcal E_t\setminus\mathcal Q}
C\bb{p_{i(P)}^\star,P}.
\label{eq:pure-death-centered-end-expansion}
\end{equation}
The same expansion holds with~$C,\kappa$ replaced by~$C',\kappa'$. If~$\mu\in\mathcal M_\star$, every centered moment in
\eqref{eq:pure-death-centered-end-expansion} is nonnegative, so
\eqref{eq:pure-death-end-patch-comparison} yields
\begin{equation*}
\mu\bb{
\prod_{P\in\mathcal E_t}
C\bb{\eta\bb{i(P)},P}
}
\leq
\mu\bb{
\prod_{P\in\mathcal E_t}
C'\bb{\eta\bb{i(P)},P}
}.
\end{equation*}
Combining this with~\eqref{eq:pure-death-bulk-patch-comparison} in the two
patch representations~\eqref{eq:patch-representation-general-law} gives
\begin{equation*}
(\mu P_t)(\chi_A)
\leq
(\mu P_t')(\chi_A),
\end{equation*}
which proves the corollary.
\end{proof}

%==============================================================================
% Section 8: discussion (now absorbed into the introduction of main.tex)
%==============================================================================

\section{Discussion}
\label{sec:discussion}

The patch representation is not restricted to patch-positive systems. Patch
positivity is one way of making the representation useful: it makes the patch
contributions nonnegative and allows termwise comparison. A different
possibility is quantitative contraction even when the contributions have mixed
signs. For example, if the relevant patch factors satisfy a uniform bound of
the form
\begin{equation*}
\abs{C(z,P)}
\leq
1-\varepsilon
\end{equation*}
for some~$\varepsilon>0$, then event where many patches exist has small contribution, exponentially in the number of patches. On the other hand,
if there are only few patches then the dual process likely died out and the patches are very long, so their terminal contributions depend only weakly on the initial configuration. This is the usual mechanism in duality arguments for ergodicity. Thus the patch representation may remain useful well
outside the patch-positive class, provided the signed patch contributions provide enough cancellations. Models such as
FA-2f, which are not patch-positive, are a natural test case for this
possibility.

The condition~$\mb p^\star\leq\mbs{\tfrac12}$ in
Corollary~\ref{cor:uniform-exponential-ergodicity} also appears to be a
feature of the proof rather than a real threshold. The value~$1/2$ enters
only because it implies that every probability measure belongs to~$\mathcal M_{-,1}$, after which Theorem~\ref{thm:common-invariant-limit}
applies directly to every initial law. Nothing in the patch representation or
the definition of patch positivity changes at this value. Corollary~\ref{cor:oriented-ergodicity} gives further evidence that the restriction is
not intrinsic: on orientable neighbourhood graphs, uniform ergodicity holds
without any assumption that~$\mb p^\star\leq\mbs{\tfrac12}$. In
particular, the perturbed East model is uniformly ergodic for every~$p\in(0,1)$ under the assumptions of Corollary~\ref{cor:perturbed-east}.

The orientability proof suggests a possible route to removing the restriction
on general neighbourhood graphs. Rather than requiring the initial law itself
to belong to~$\mathcal M_\star$ or~$\mathcal M_-$, one may try to prove
that after a burn-in period its evolution approaches one of these classes on
finite sets. A natural target is
\begin{equation*}
\liminf_{t\to\infty}
\inf_{\eta\in\Omega}
P_t\chi_A^\star(\eta)
\geq0
\qquad
\text{for every }A\Subset\Lambda,
\end{equation*}
or an analogous approximation to~$\mathcal M_-$. Once such a statement is
available, Theorem~\ref{thm:common-invariant-limit} gives convergence to the invariant measure unique in~$\mathcal{M}_-$. Orientability provides one mechanism for this
burn-in by reducing the problem locally to a finite system with autonomous
boundary. A replacement for this mechanism on general neighbourhood
graphs would extend the convergence result beyond the present restriction on
the initial law.

%==============================================================================
% Parked from main.tex (Section 4.3, 2026-10-06): Bernoulli product law form
% of the patch representation. Not used in the ergodicity paper.
%==============================================================================

For a Bernoulli product measure~$\mu_{\mb p}$ with
$\mu_{\mb p}\bb{\eta(i)=1}=p_i$, the sites~$i(P)$,
$P\in\mathcal E_t$, are distinct. Hence independence of the spins at these
sites and~$z\mapsto C(z,P)$ being affine gives
\begin{equation*}
\mu_{\mb p}\bb{
\prod_{P\in\mathcal E_t}C\bb{\eta\bb{i(P)},P}
}
=
\prod_{P\in\mathcal E_t}
C\bb{p_{i(P)},P}.
\end{equation*}
Therefore
\begin{equation}
\label{eq:patch-representation-product-law}
(\mu_{\mb p}P_t)(\chi_A)
=
\E_A\Cb{
\prod_{P\in\mathcal B_t}C(P)
\prod_{P\in\mathcal E_t}C\bb{p_{i(P)},P}
}.
\end{equation}

